\chapter[Referencial Teórico]{Referencial Teórico}
\label{cap:referencial}

Neste capítulo serão abordados os principais conceitos e assuntos associados aos estudos e as realizações importantes para a melhor compreensão de seu contexto e o desenvolvimento condizente as demandas identificadas.

\section[Tutoria]{Tutoria}
\label{sec:tutoria}

A prática de uma pessoa mais experiente ensinar ou orientar uma pessoa mais novata no estudo ou no exercício de uma determinada atividade pode receber diversos nomes a depender do contexto: mentoria, tutoria ou monitoria.

Sendo assim, é importante explicar cada um dos conceitos citados para melhor discernimento do que o projeto visa alcançar, contribuindo ainda com o melhor esclarecimento sobre cada um deles e suas diferenças.

Mentor é um indivíduo experiente que tem por objetivo passar adiante suas experiências e vivências em uma determinada área de atuação. Normalmente age como uma referência, uma pessoa disposta a demonstrar tudo o que sabe, orientando, apontando direções e estabelecendo desafios que permitam ao aprendiz que está sendo orientado, denominado mentorado, desenvolver seu conhecimento a partir de suas experiências \cite{oliveira2020mentoria}.

A mentoria é um acompanhamento de longo prazo e marcada por uma relação de respeito e confiança entre o mentor e o mentorado, envolvendo aspectos de carreira e apoio psicológico \cite{oliveira2020mentoria}. Entretanto, o mentor tem um papel que ultrapassa a orientação para estudo, buscando também o desenvolvimento pessoal e psicossocial do mentorado \cite{botti2008preceptor}. Exemplos incluem a aquisição de competências socioemocionais e cognitivas e a melhoria da autoestima e autorregulação \cite{opp2024mentoria}.

Essa natureza pessoal da mentoria é presente desde o poema que deu origem a essa palavra: \textit{Odisseia} de Homero, em que Odisseu confia a Mentor, seu fiel e sábio amigo, a educação de seu filho Telêmaco antes de partir para a guerra de Tróia. Atena, deusa da sabedoria e da razão, assume a forma de Mentor e orienta Telêmaco no seu desenvolvimento físico, intelectual e espiritual \cite{semiao2009tutoria}.

Tutoria pode ser definida como uma relação de ajuda entre duas pessoas com diferentes níveis de conhecimento, competências e/ou experiências em que o tutorando, o indivíduo de menor capacidade, recebe apoio do tutor para superar dificuldades escolares ou acadêmicas \cite{opp2024mentoria}. O tutor não necessariamente precisa ser um docente formal, podendo ser um familiar, estudante ou um voluntário \cite{topping2000tutoring}.

A tutoria existe desde a Antiguidade, onde a cultura de povos e valores como honra, justiça e honestidade eram passados pelos tutores para seus aprendizes. A partir das universidades na Idade Média a tutoria passou também a ser acadêmica, com tutores sendo alunos mais velhos que ensinavam habilidades como leitura e interpretação de texto para os estudantes mais jovens, sob supervisão de um docente \cite{geib2007tutoria}.

Comparado com o aprendizado tradicional com um professor, a tutoria pode fornecer mais experiência prática, mais atividades e variedade e também mais ajuda individualizada \cite{topping2000tutoring}. Embora parecidas, a tutoria e a mentoria se diferem em seus objetivos, visto que a tutoria visa ajudar na superação de dificuldades através do ensino de conteúdos ou competências específicas, sem buscar ativamente o desenvolvimento pessoal \cite{opp2024mentoria}. Entretanto, este ainda pode vir a ocorrer como consequência das interações, como aumento da autoconfiança e desenvolvimento de pensamento crítico, tanto para o tutor quanto para o tutorando \cite{lin2025peer}.

Monitoria é uma atividade que surgiu na Idade Média, e consiste em um estudante mais experiente de uma turma auxiliar o professor no ensino de colegas de classe mais novatos. Na metade do século XIV, os professores tinham quase sempre um monitor que os auxiliavam na escolarização. No século XVIII existia o Método Monitorial de Lancaster na Inglaterra, em que adolescentes, instruídos diretamente pelos mestres, atuavam como monitores ensinando outros adolescentes \cite{frison2016monitoria}. Nota-se que a monitoria continua a ser usada até os dias atuais.

No Brasil, o programa de monitoria está regulamentado pelas instituições de ensino superior e tem como finalidade promover a participação de estudantes com desempenho destacado em disciplinas, atuando como facilitadores da aprendizagem dos colegas e colaborando com o docente na construção de metodologias de ensino \cite{monitoria2025}.

A teoria por trás da monitoria tem origem no pensamento pedagógico de Johann Heinrich Pestalozzi (1746 a 1827), que rejeitava a educação mecanicista e defendia o desenvolvimento humano integral. Para Pestalozzi (\citeyear{pestalozzi1946monitoria}), a mediação entre estudantes tem valor pedagógico justamente porque não segue um modelo de produção em série. Trata-se de uma atividade colaborativa entre aprendizes, baseada na cooperação e no afeto. Ao trazer essa ideia para a monitoria acadêmica, fica claro que a melhor forma de aprender é ensinando \cite{aranha2006monitoria}.

Assim como na tutoria, ambos monitor e estudante se beneficiam da monitoria. O estudante recebe apoio no aprendizado do conteúdo da disciplina e o monitor desenvolve competências acadêmicas, sociais e profissionais, como comunicação, liderança, empatia e didática, além de seu engajamento em programas de monitoria poder abrir portas para novas oportunidades acadêmicas, como a participação em projetos de pesquisa e extensão \cite{monitoria2025}.

Diante do exposto, o modelo de interação pretendido pelo projeto é a tutoria, focada no aprendizado de conteúdos e habilidades específicas, não sendo o objetivo o desenvolvimento pessoal em conjunto com o aprendizado, inerente à mentoria, e não possuindo o caráter acadêmico intrínseco à monitoria. Entretanto, a tutoria pode evoluir para uma mentoria caso se forme um vínculo entre o tutor e o tutorado, em que e o tutor queira ajudar nos problemas pessoais do tutorado, passando a ser um mentor para o mesmo.

Neste trabalho, a aprendizagem é compreendida, principalmente, sob a ótica da Teoria da Aprendizagem Significativa de David Paul Ausubel, que define aprendizagem como a organização e integração de informações na estrutura cognitiva do indivíduo. De acordo com essa teoria, a aprendizagem é significativa quando os novos conceitos são compreendidos e relacionados com conhecimentos prévios do indivíduo que está aprendendo, criando conexões significativas e sendo mais provável de ser transferida para a memória de longo prazo, enquanto que uma aprendizagem mecânica ocorre quando o indivíduo não realiza tais conexões, simplesmente memorizando fatos e informações sem entender os princípios ou conexões subjacentes, podendo ocorrer descarte congnitivo de tal conhecimento sem relevância, sem significado e sem sentido \cite{junior2023olhar}.

Considerando que o fator isolado mais importante para a aprendizagem seria o que o aprendiz já sabe, a tutoria se apresenta como um mecanismo ideal para esse processo, pois permite que o tutor identifique os conhecimentos prévios do aprendiz e facilite a assimilação da nova informação, promovendo a retenção do conhecimento e evitando a aprendizagem mecânica.

Adicionalmente, a dinâmica da tutoria é compreendida neste trabalho como um processo de aprendizagem colaborativa, fundamentado na perspectiva sociointeracionista de Lev Vygotsky. Segundo \citeonline{saul2025}, essa abordagem concebe o desenvolvimento cognitivo e o aprendizado como inerentemente sociais, mediados pela interação entre indivíduos com diferentes níveis de conhecimento. Nesse cenário, o aprendizado ocorre quando um indivíduo realiza uma tarefa ou atividade que está além de suas capacidades recebendo o apoio de uma pessoa ou recurso com mais conhecimento do que ele, denominado More Knowledgeable Other (MKO, ou "Outro Mais Conhecedor"). Através do compartilhamento do conhecimento, dicas e encorajamento, o aprendiz acessa sua Zona de Desenvolvimento Proximal (ZDP), justamente o espaço entre o que consegue fazer sozinho e com ajuda, internalizando novas habilidades e podendo se tornar um MKO para outros aprendizes.

A integração de recursos tecnológicos no suporte ao aprendizado encontra forte sustentação na convergência de abordagens pedagógicas clássicas e contemporâneas. Conforme preconizado por Rissoli, Giraffa e Martins (\citeyear{rissoli2006}), a modelagem de Sistemas Tutores Inteligentes (STI) permite aplicar com sucesso os conceitos da Teoria da Aprendizagem Significativa de Ausubel, utilizando a estrutura lógica para identificar os subsunçores do aluno e recomendar conteúdos personalizados. Essa ancoragem cognitiva, quando associada aos princípios de mediação ZPD de Vygotsky, fornece o arcabouço ideal para sistemas de recomendação que guiam o estudante em seu ritmo próprio. Adicionalmente, a transposição desses conceitos para ambientes virtuais modernos resgata a visão de Pestalozzi sobre a importância do acompanhamento individualizado e do aspecto afetivo/motivacional no ensino, o qual, no escopo deste trabalho, é potencializado por meio de dinâmicas de gamificação e algoritmos de recomendação web.

Para viabilizar a busca e o contato com os tutores de forma acessível, a solução proposta é o desenvolvimento de uma aplicação web. Esta aplicação servirá como um ambiente digital onde os usuários poderão não apenas buscar por especialistas em determinados conteúdos ou habilidades, mas também realizar o contato inicial e agendamento das sessões ou encontros que propiciarão aprendizado(s).

\section[Aplicação Web]{Aplicação Web}
\label{sec:aw}

% O projeto visa desenvolver uma aplicação web.
Uma aplicação web é um programa computacional que é executado através de um navegador de Internet, também chamado de \textit{browser}, como Google Chrome, Mozilla Firefox e outros, com o objetivo de executar funções, como um formulário de contato ou uma calculadora. Exemplos de aplicações web são Google Docs e Microsoft Office 365 \cite{harwood2022internet}.

Segundo \citeonline{shklar2009web}, uma aplicação web é diferente de um \textit{website}. Um \textit{website} é essencialmente informativo, projetado para fornecer conteúdo de forma estática, semelhante a um folheto digital ou a um \textit{site} de notícias, não necessariamente oferecendo funcionalidades. Essa distinção é reforçada por \citeonline{delaRosa2025webapp}, que destaca diferenças adicionais, como o nível de interatividade. \textit{Websites} tendem a oferecer menor interatividade, geralmente limitando-se a ações básicas, como navegação por \textit{links}, enquanto aplicações web permitem interações mais complexas, como envio de arquivos, preenchimento de formulários e edição de documentos, por serem projetados para oferecer funcionalidades para o usuário. 

Aplicações web são programas elaborados sob uma arquitetura cliente/servidor \cite{shklar2009web}. Segundo \citeonline{nyabuto2024architectural}, essa arquitetura funciona baseada no envio de requisições e respostas, em que o cliente inicia a comunicação enviando um pedido (requisição) e o servidor processa esse pedido e retorna uma resposta ao cliente. Na aplicação proposta, um exemplo de comunicação é o aprendiz buscar tutores, enviando uma requisição para o servidor, que a processa e retorna uma lista de tutores (resposta) para o aprendiz.

A interação do usuário com uma aplicação web pode ser considerada um diálogo. O cliente (usuário) fornece dados ou clica em elementos da interface, como um botão, para realizar uma ação. A aplicação então fornece uma resposta para a ação do usuário, como exibir dados em resposta a submissão de um formulário de pesquisa, que pode então gerar uma nova ação do usuário \cite{harwood2022internet}.

Diferente das aplicações para computador tradicionais, que exigem instalação e configuração individual em cada máquina, as aplicações web residem em servidores remotos, o que altera fundamentalmente a forma como o software é distribuído e mantido. Segundo \citeonline{randy2022}, essa característica elimina a barreira da instalação local e simplifica bastante o ciclo de atualização: uma vez que o código é atualizado no servidor, todos os usuários passam a utilizar a versão mais recente de forma instantânea e simultânea quando acessam os serviços oferecidos por este servidor.

Além disso, ainda segundo \citeonline{randy2022}, as aplicações web oferecem acessibilidade universal, por serem independentes do sistema operacional: desde que o dispositivo possua um navegador moderno e compatível, a aplicação funcionará independentemente de o usuário estar em um ambiente Windows, macOS, Linux ou outros.

Para enriquecer a experiência na aplicação (software), serão implementadas duas estratégias principais. Primeiramente, um sistema de recomendação será desenvolvido para garantir que os usuários encontrem tutores compatíveis com seus interesses e necessidades. Adicionalmente, para promover o uso contínuo e manter os usuários engajados, serão aplicados elementos de Gamificação.

\subsection[Sistema de Recomendação]{Sistema de Recomendação}
\label{sec:recomendação}

Os Sistemas de Recomendação (SR) são técnicas e ferramentas de software que fornecem sugestões de itens (termo geral usado para denotar o que o sistema está recomendando) com grande probabilidade de serem relevantes para um usuário em particular \cite{springerSR2022}.

Normalmente os SR estão em aplicações web, aprendendo com interações anteriores dos usuários e/ou com atributos dos itens \cite{Aggarwal2016}. Esses sistemas ajudam os usuários a encontrarem itens de interesse em grandes conjuntos de dados, priorizando ou filtrando alternativas com base no histórico do usuário, preferências ou comportamentos similares de outros usuários \cite{springerSR2022}.

Segundo \citeonline{patel2023recommendation}, os Sistemas de Recomendação Baseados em Conteúdo (SRBC) recomendam itens com base nos dados que um usuário fornece direta ou indiretamente sobre os itens que avaliou ou que, no caso de lojas digitais, comprou. É realizada uma análise das descrições de itens previamente avaliados pelo usuário. Em seguida, um perfil de usuário é criado com base nas tendências de classificação anteriores, sendo então utilizado para gerar recomendações de documentos, artigos, notícias, filmes, entre outros. A premissa é que o usuário terá interesse em itens similares àqueles com os quais já interagiu positivamente no passado.

Assumindo um cenário de loja digital, o SRBC funciona da seguinte maneira: o sistema identifica um item que foi avaliado positivamente ou comprado pelo usuário no passado. É analisado o conteúdo, a categoria e as características desse item para descobrir quais outros itens possuem atributos semelhantes. Com base nas pontuações de similaridade, o sistema recomenda o item mais relevante ao usuário \cite{patel2023recommendation}.

O Sistema de Recomendação Colaborativo (SRC) recomenda itens baseando-se nas recomendações de diferentes pessoas. A ideia central por trás dessa técnica é que pessoas semelhantes compartilham os mesmos interesses e, portanto, o usuário gostará de itens semelhantes \cite{kumar2021recommender}.

A abordagem colaborativa usa uma matriz de classificações de itens de usuários como única entrada e normalmente produzem os seguintes tipos de saída: uma previsão (numérica) indicando até que ponto o usuário atual gostará ou não gostará de um determinado item e uma lista de \textbf{\textit{n}} itens recomendados. Essa lista dos \textbf{\textit{n}} melhores itens não deve, obviamente, conter itens que o usuário atual já tenha adquirido \cite{Jannach2011}. 

Segundo \citeonline{kumar2021recommender}, como vantagens do SRC pode-se citar sua independência do conteúdo, visto que não precisa analisar as características intrínsecas do item (como gênero ou descrição técnica), dependendo apenas das avaliações dos usuários, e também a possibilidade de surpreender o usuário, pois é capaz de recomendar itens que o usuário não esperaria ou que são diferentes do seu perfil habitual, desde que usuários similares tenham gostado desses itens. Como desvantagens, é citada a dificuldade de recomendação para usuários novos, pois não há histórico de comportamento disponível para gerar sugestões precisas, e a esparsidade de dados: como a maioria dos usuários avalia apenas uma pequena fração dos itens disponíveis, a matriz de interação fica com muitos espaços vazios, o que pode reduzir o desempenho do sistema.

Os Sistemas de Recomendação Híbridos (SRH) surgem como uma opção interessante ao combinar diversas técnicas ou algoritmos, visando unir os pontos positivos de diferentes métodos, a fim de atender suas fraquezas e melhorar a eficiência das recomendações. O SRC é baseado na avaliação das preferências dos usuários e no comportamento coletivo da comunidade, a fim de produzir sugestões com base nas semelhanças entre perfis de usuários. Em contrapartida, o SRBC utiliza de atributos próprios dos itens, como suas características estruturais e descrições em texto, para determinar a relevância pessoal de interações. Embora ambas as metodologias tenham benefícios reconhecidos, cada uma tem suas limitações específicas — como a limitação de dados na filtragem colaborativa quando pouco ou nenhum perfil de usuário para comparação e a insuficiência de variedade nas recomendações que se baseiam apenas no conteúdo \cite{Jannach2011}.

Contudo, variações na arquitetura dos sistemas de recomendação também são utilizados no ambiente educacional. Por exemplo, no Sistema de Apoio Educacional (SAE), o módulo conhecido como Material, foi integrado a este software educacional que corresponde a um ITA (Inteligent Teaching Assistant), a fim de fornecer recomendações de conteúdo de apoio educacional conforme estejam sendo os resultados do desempenho em relação ao esforço e ao empenho de cada estudante usuário do SAE.
Este módulo Material, que corresponde a um sistema de recomendação de conteúdo e de pessoas, foi integrado na arquitetura ITA que acompanha personalizadaemente cada estudante usuário do SAE e fornece recomendações de conteúdos de estudo para cada demanda que o estudante o solicite. No entanto, na proposta deste módulo do SAE (Material) também são fornecidas as recomendações de colegas de turma que estão tendo um melhor desempenho que o estudante usuário do SAE, a fim de contribuir com o possível estudo colaborativo entre integrantes de uma mesma turma que possuem resultados diferentes e acompanhados em tempo real por este software web. Assim, o módulo Material identifica quem está em melhores condições de colaborar com a aprendizagem e o recomenda ao estudante usuário do SAE \cite{tcc-viana-e-bueno}.

% [TCC Ênio Carlos Viana e Willian Fernandes de Oliveira Bueno. Sistema de Recomendações. 2014. Trabalho de Conclusão de Curso. (Graduação em Ciência da Computação) - Universidade Católica de Brasília. Orientador: Vandor Roberto Vilardi Rissoli.]

Este software educacional ainda fornece recomendação de quais estudantes mais experientes (monitor) naquele conteúdo que o usuário do SAE está solicitando auxílio seu auxílio através de recomandação de conteúdos de estudo. Este sistema averigua qual monitor obteve melhor desempenho naquele conteúdo em específico quando também era aluno, procurando assim contribuir com o processo de aprendizagem de todos os envolvidos estimulando o estudo colaborativo para a assimilação ser mais significativa \cite{simposio2013}.

% [artigo USP CAMPOS, B. S. ; VIANA, E. C. ; MARTINELLO, F. X. ; RISSOLI, V. R. V. . A atuação do Sistema de Recomendação no processo de ensino-aprendizagem. In: Simpósio Internacional de Iniciação Científica, 2013, São Carlos, SP. Anais do 21 Simpósio Internacional de Iniciação Científica da Universidade de São Paulo. São Carlos: USP, 2013. v. 21.]                                                                                                                          ou                                                                                                                        [VIANA, E. C. ; BUENO, W. F. O. ; RISSOLI, V. R. V. . Sistema de Recomendações de Conteúdo para Apoio à Aprendizagem Significativa. In: Congresso de Iniciação Científica da Universidade de Ribeirão Preto, 2012, Ribeirão Preto. Anais do 13 Congresso de Iniciação Científica da Universidade de Ribeirão Preto. Ribeirão Preto: UNAERP, 2012. v. 13.]

\subsection[Gamificação]{Gamificação}
\label{sec:gamificacao}

A gamificação é a infusão intencional de elementos e mecânicas de jogos em situações que não são de jogos. O objetivo principal é engajar os participantes ativamente, motivá-los e ajudá-los a alcançar objetivos específicos. Essa abordagem aproveita nossas tendências naturais, como o desejo de competição, o senso de realização, o trabalho em equipe e o altruísmo, para impulsionar a participação, o desempenho e a satisfação geral. A gamificação envolve várias estratégias, como a introdução de sistemas de pontuação, a concessão de emblemas (\textit{badges}) de conquista e a criação de tabelas de classificação (\textit{leaderboards}). Esses elementos adicionam uma dimensão interativa e recompensadora as atividades que não são inerentemente jogos, como educação, \textit{marketing} e tarefas no local de trabalho \cite{deliyannis2023gamification}.

Pesquisas interdisciplinares demonstram que os jogos são mediadores significativos entre a aprendizagem e a socialização ao longo da vida, fazendo com que a Aprendizagem Baseada em Jogos (GBL) se tornasse um importante ponto de diálogo na educação. A experiência de jogar oferece múltiplas oportunidades para aprender e interagir com o ambiente \cite{magtani2024design}

Como uma das metodologias ativas para aprendizagem, a gamificação se apresenta como técnica capaz de impulsionar o envolvimento nas atividades educacionais. Gamificar atividades educacionais não significa trazer ou realizar jogos em sala de aula. Antes é incorporar os conceitos, dinâmicas, interação e desafios presentes no universo dos jogos para potencializar o processo de aprendizagem \cite{souza2025metodologias}.

Para compreender como os elementos de jogos influenciam o comportamento do usuário, \citeonline{chou2015actionable} propõe o \textit{framework Octalysis}. Diferente de abordagens que focam apenas na mecânica superficial, esta estrutura analisa a gamificação sob a ótica da motivação humana, identificando oito \textit{Core Drives} (Impulsos Centrais) que orientam nossas ações e decisões:

\begin{itemize}
    \item \textbf{Significado Épico e Chamado (\textit{Epic Meaning \& Calling}):} Ocorre quando o usuário acredita estar engajado em algo maior que ele mesmo ou que foi ``escolhido'' para uma missão. É o pilar que motiva, por exemplo, contribuidores de projetos colaborativos como a Wikipédia, que dedicam tempo voluntário por uma causa maior.

    \item \textbf{Desenvolvimento e Realização (\textit{Development \& Accomplishment}):} Trata-se do impulso interno de progredir, desenvolver habilidades e superar desafios. Este é o motor principal por trás da maioria dos sistemas de pontos, medalhas e rankings, pois os usuários sentem satisfação ao visualizar seu crescimento e completar tarefas difíceis.

    \item \textbf{Empoderamento da Criatividade e Feedback (\textit{Empowerment of Creativity \& Feedback}):} Envolve engajar o usuário em um processo criativo em que precisa descobrir novas combinações e estratégias. O engajamento surge da autonomia para testar diferentes abordagens e receber \textit{feedback} imediato sobre o resultado de suas escolhas.

    \item \textbf{Propriedade e Posse (\textit{Ownership \& Possession}):} Baseia-se na motivação gerada quando o usuário sente que possui algo. Ao sentir-se dono de um item (virtual ou não), o usuário tende a querer melhorá-lo, protegê-lo e acumular mais. Este impulso justifica o uso de bens virtuais, coleções e personalização de perfis.

    \item \textbf{Influência Social e Relacionamento (\textit{Social Influence \& Relatedness}):} Incorpora todos os elementos sociais que motivam as pessoas, como mentoria, aceitação social, competição e companheirismo. A observação do comportamento de outros usuários e o desejo de pertencimento são fortes catalisadores de ação.

    \item \textbf{Escassez e Impaciência (\textit{Scarcity \& Impatience}):} É o impulso de desejar algo simplesmente porque é raro, exclusivo ou difícil de obter imediatamente. A limitação de tempo ou de acesso cria um senso de urgência que pode acelerar a tomada de decisão do usuário.

    \item \textbf{Imprevisibilidade e Curiosidade (\textit{Unpredictability \& Curiosity}):} Relaciona-se ao engajamento gerado pela incerteza do que virá a seguir. A curiosidade cognitiva motiva o usuário a continuar interagindo com o sistema para descobrir novos resultados ou conteúdos.

    \item \textbf{Perda e Aversão (\textit{Loss \& Avoidance}):} É a motivação baseada no medo de perder algo conquistado ou de que um evento negativo ocorra. No contexto de gamificação, pode ser utilizado para incentivar a continuidade do uso, evitando que o progresso ou status acumulado seja perdido.
\end{itemize}

A aplicação desses impulsionadores permite que sistemas gamificados, como a aplicação proposta neste trabalho, utilizem elementos como pontos e conquistas não apenas como recompensas arbitrárias, mas como ferramentas estruturadas para satisfazer necessidades psicológicas de competência e progresso.

Um sistema gamificado é composto de elementos que podem ser encontrados em \textit{games}. Alguns destes elementos podem ser definidos como Pontos, Níveis, \textit{Rankings}, Desafios/Missões e Conquistas \cite{klock2014gamificacaoAVA}.

Os pontos são abertos, diretos e motivacionais, permitindo a utilização de vários tipos diferentes de pontuação, dependendo do objetivo. Alguns dos tipos de pontos podem ser descritos como pontos de experiência, pontos de habilidade e pontos de reputação \cite{zichermann2011gamificacao}.

Os níveis podem indicar o progresso do usuário dentro do sistema gamificado. Os níveis podem ser definidos como níveis de jogo, níveis de dificuldade e níveis do usuário \cite{zichermann2011gamificacao}.

Os \textit{Rankings} são a comparação entre usuários envolvidos, a fim de visualizar suas progressões dentro do ambiente e inserir um senso de competição entre eles \cite{klock2014gamificacaoAVA}. Os desafios e missões norteiam os usuários sobre as atividades que devem ser feitas dentro do sistema \cite{fadel2014gamificacao}.

Conquistas são uma outra forma de pontos, possuindo uma representação visual de alguma realização do usuário no sistema \cite{hunter2012gamificacao}.

A proposta, que será apresentada mais detalhadamente no Capítulo 3, é desenvolver uma aplicação web (programa computacional executado através do navegador de Internet) para facilitar o contato entre um tutor, que busca ajudar alguém a aprender ou se aperfeiçoar em uma habilidade, e um aprendiz que busca auxílio em algo que deseja ou precisa aprender. A aplicação contará com um sistema de recomendação para sugerir tutores relevantes a cada usuário e gamificação (elementos de jogos digitais, como pontos e níveis) para manter tanto tutores quanto aprendizes engajados com o uso dessa aplicação.